\documentclass[../../main.tex]{subfiles}
\begin{document}

\begin{flushright}
\footnotesize\textit{【自動文字起こし・要確認】}
\end{flushright}

[解] 高さ $y$ の時の表面の面積 $S$、円錐の半径を $r$ とすると
\[ 10 - y = S \dfrac{dy}{dt} \cdots \text{①} \]
$\therefore S = \pi \left( \dfrac{r(10-y)}{10} \right)^2$ を代入してセイリ
\[ dt = \dfrac{\pi r^2}{100}(10-y)dy \]
両辺積分して
\[ t = \dfrac{\pi r^2}{100}\left(10y - \dfrac{1}{2}y^2\right) + C \]
$t=0$ で $y=0$ から $C=0$。$(t, y) = (540, 2), y=10$ を各々代入
\[ \begin{cases} 540 = \dfrac{\pi r^2}{100} \cdot 18 \\ t = \dfrac{\pi r^2}{100} \cdot 50 \end{cases} \]
辺々わって
\[ \dfrac{t}{540} = \dfrac{50}{18} \cdot 54 \]
\[ \therefore t = 1500 \]
だから、タンクがいっぱいになるのは、あと
\[ \dfrac{1500}{60} - 9 \text{ (h)} = 16 \text{ (h)} \]
後である。

\end{document}
